\section{Related Work}

\paragraph{LLM-as-a-judge reliability.} Using a language model to score another
model's output is itself a measurement with known failure modes, and our result
is a direct addition to this literature. \citet{shi2024judging} find systematic
(non-random) position bias across 15 LLM judges and $>150{,}000$ evaluation
instances, motivating our use of blinded, randomly ordered single-response
scoring rather than pairwise comparison. \citet{wataoka2024selfpreference} find
that GPT-4 exhibits measurable self-preference bias, and
\citet{pombal2026rubricspb} extend this to rubric-based (not just pairwise)
evaluation, showing self-preference bias persists under objective criteria and
that ensembling multiple judges mitigates but does not eliminate it. Prior work
largely quantifies judge \emph{bias} as a deviation from a reference or a
preference inflation; we contribute a case in which the choice of judge does not
merely shift scores but \emph{inverts the substantive behavioral conclusion} of
a preregistered evaluation, and in which the largest disagreement coincides with
a judge scoring its own model family. We frame this as an observed
measurement-stability failure, not as a demonstration of a specific causal bias
mechanism.

\paragraph{Goal conflict and agentic misalignment.} Our behavioral probe is
motivated by \citet{anthropic2025agentic}'s finding that goal conflict with a
deployer, or a perceived threat to a model's continued operation, can
independently elicit harmful autonomous action across frontier and mid-size
models. We study a lower-stakes textual precursor --- persuasive-communication
behavior under an analogous goal-pressure framing, in small open-weight models
--- but our conclusion is about how reliably such behavior can be
\emph{measured}, not about the behavior's prevalence.

\paragraph{Manipulative design and dark patterns.} The Center for Democracy and
Technology's taxonomy identifies 37 manipulative design patterns across five
categories in AI chatbots \citep{cdt2026darkpatterns}. Our six rubric dimensions
were designed independently but correspond to several of these categories (e.g.
deceptive\_omission to informationally misleading design,
respect\_for\_refusal/transparency to autonomy-compromised-for-engagement),
providing an external cross-check that the rubric targets recognized real-world
manipulation categories rather than an ad hoc list.

\paragraph{Evaluation awareness.} \citet{needham2025evalaware} find that frontier
models substantially exceed chance at detecting whether they are in an
evaluation context (Gemini-2.5-Pro AUC 0.83, below the human baseline of 0.92).
This motivates measuring, rather than assuming away, evaluation-awareness
signals; our rubric includes a diagnostic (non-composite)
\texttt{evaluation\_awareness} dimension, though we do not separately test our
small subject models' evaluation-detection capability.

\paragraph{Statistical methodology for repeated-measures behavioral data.}
Because each scenario is evaluated under all four conditions by each subject
model, observations are not independent. \citet{mackinnon2023clusterboot} provide
the methodological basis for the scenario-blocked bootstrap and
scenario-clustered inference we adopt in place of treating 180 outputs as 180
independent draws.
